\documentclass[11pt,a4paper]{article}
\usepackage{fontspec,amsmath,amssymb,amsthm,microtype}
\setmainfont{lmroman10-regular.otf}[BoldFont=lmroman10-bold.otf,ItalicFont=lmroman10-italic.otf,BoldItalicFont=lmroman10-bolditalic.otf]
\usepackage{xcolor}
\usepackage[most]{tcolorbox}
\usepackage{mathrsfs}
\providecommand{\lt}{<}
\providecommand{\gt}{>}
\setmonofont{DejaVuSansMono.ttf}[BoldFont=DejaVuSansMono-Bold.ttf,ItalicFont=DejaVuSansMono-Oblique.ttf,BoldItalicFont=DejaVuSansMono-BoldOblique.ttf,Scale=0.86]
\newfontfamily\leanfallback{STIXTwoMath-Regular.otf}[Scale=0.86]
\definecolor{leanrule}{gray}{0.62}
\definecolor{leanback}{gray}{0.962}
\definecolor{leantitle}{gray}{0.22}
\definecolor{leanmuted}{gray}{0.38}
\newtcolorbox{leanbox}[1]{enhanced,breakable,colback=leanback,colframe=leanrule,
  boxrule=0.4pt,arc=0.6mm,left=2.4mm,right=2.4mm,top=1.8mm,bottom=1.8mm,
  before skip=0.8em,after skip=1em,fonttitle=\small\bfseries,coltitle=leantitle,
  colbacktitle=leanback,titlerule=0.3pt,title={#1}}
\newlength{\leancharwidth}
\newenvironment{leancode}{\par\smallskip\begingroup\ttfamily\footnotesize\raggedright
  \settowidth{\leancharwidth}{M}\setlength{\parskip}{0pt}\setlength{\parindent}{0pt}
  \frenchspacing}{\par\endgroup\smallskip}
\newcommand{\leanline}[2]{\hangindent=\dimexpr#1\leancharwidth+4\leancharwidth\relax
  \hangafter=1\noindent\hspace*{#1\leancharwidth}#2\par}
\newcommand{\leansource}[1]{\par{\footnotesize\color{leanmuted}#1}\par}
\newcommand{\leanlabel}[1]{\par\smallskip\noindent{\small\bfseries #1}\par}
\usepackage[margin=25mm]{geometry}
\usepackage[unicode,hidelinks]{hyperref}
\hypersetup{pdfauthor={LangevinLean Project}}
\setlength{\emergencystretch}{3em}
\setlength{\parskip}{0.4em}
\setlength{\parindent}{0pt}
\allowdisplaybreaks[2]
\raggedbottom
\title{\texorpdfstring{Actual paired independent innovation paths}{Actual paired independent innovation paths}}
\author{LangevinLean Project}
\date{}
\begin{document}\maketitle
Draft mathematical exposition; these notes may contain mistakes.


Two independent paths of independent identically distributed innovations can be zipped coordinatewise into a single path of independent pairs. The exact identity persists when an independent initial state is present.


\section{\texorpdfstring{Zipping the two actual infinite product laws}{Zipping the two actual infinite product laws}}


\textbf{Claim.} If \((\xi_n)_n\) and \((U_n)_n\) are independent paths with laws \(\mu^{\otimes\mathbb N}\) and \(\nu^{\otimes\mathbb N}\), their zipped path has law \((\mu\otimes\nu)^{\otimes\mathbb N}\).



\textbf{Proof.} Restrict to any finite set of times. The two restricted laws are finite product laws. Reordering their coordinates gives the finite product of pair laws. The infinite product measure is determined by these finite coordinate cylinders, proving the full path identity.

\begin{leanbox}{Lean: \texttt{paired\_\allowbreak{}infinitePi}}
\leansource{\texttt{paired\_\allowbreak{}infinitePi} \quad WassersteinMixing/\allowbreak{}ProductNoise.lean, lines 17–54}
\begin{leancode}
\leanline{0}{theorem\ paired\_infinitePi\ (μ\ :\ Measure\ A)\ (ν\ :\ Measure\ B)}
\leanline{4}{[IsProbabilityMeasure\ μ]\ [IsProbabilityMeasure\ ν]\ :}
\leanline{4}{((Measure.infinitePi\ fun\ \_\ :\ ℕ\ ↦\ μ).prod}
\leanline{6}{(Measure.infinitePi\ fun\ \_\ :\ ℕ\ ↦\ ν)).map}
\leanline{8}{(fun\ p\ n\ ↦\ (p.1\ n,\ p.2\ n))\ =}
\leanline{6}{Measure.infinitePi\ (fun\ \_\ :\ ℕ\ ↦\ μ.prod\ ν)}
\end{leancode}
\leanlabel{How the Lean proof follows the argument}
\texttt{paired\_\allowbreak{}infinite\allowbreak{}Pi} first identifies the finite restriction pushforwards and then applies equality of infinite product measures from their cylinder masses. The proof uses actual measure-preserving coordinate reorderings.
\end{leanbox}

\section{\texorpdfstring{Appending a fresh decision path to an actual source}{Appending a fresh decision path to an actual source}}


\textbf{Claim.} If an existing source has joint law \(\mu\otimes\eta^{\otimes\mathbb N}\) for its initial state and innovation path, adjoining an independent path with law \(\nu^{\otimes\mathbb N}\) gives initial/\allowbreak{}paired-path law \(\mu\otimes(\eta\otimes\nu)^{\otimes\mathbb N}\).



\textbf{Proof.} Take the actual product extension. The original joint pushforward is preserved on its first coordinate. Reassociate the three factors, apply the zipped-path identity to the last two, and map the result back to the initial/\allowbreak{}paired-path representation.

\begin{leanbox}{Lean: \texttt{augment\_\allowbreak{}initial\_\allowbreak{}noise\_\allowbreak{}law}}
\leansource{\texttt{augment\_\allowbreak{}initial\_\allowbreak{}noise\_\allowbreak{}law} \quad WassersteinMixing/\allowbreak{}ProductNoise.lean, lines 65–102}
\begin{leancode}
\leanline{0}{theorem\ augment\_initial\_noise\_law\ (Q\ :\ Measure\ Ω)\ [IsProbabilityMeasure\ Q]}
\leanline{4}{(μ\ :\ Measure\ X)\ (η\ :\ Measure\ A)\ (ν\ :\ Measure\ B)}
\leanline{4}{[IsProbabilityMeasure\ μ]\ [IsProbabilityMeasure\ η]\ [IsProbabilityMeasure\ ν]}
\leanline{4}{(initial\ :\ Ω\ →\ X)\ (noise\ :\ Ω\ →\ ℕ\ →\ A)}
\leanline{4}{(hi\ :\ AEMeasurable\ initial\ Q)\ (hn\ :\ AEMeasurable\ noise\ Q)}
\leanline{4}{(hjoint\ :\ Q.map\ (fun\ ω\ ↦\ (initial\ ω,\ noise\ ω))\ =}
\leanline{6}{μ.prod\ (Measure.infinitePi\ fun\ \_\ :\ ℕ\ ↦\ η))\ :}
\leanline{4}{(Q.prod\ (Measure.infinitePi\ fun\ \_\ :\ ℕ\ ↦\ ν)).map}
\leanline{6}{(fun\ p\ ↦\ (initial\ p.1,\ fun\ n\ ↦\ (noise\ p.1\ n,\ p.2\ n)))\ =}
\leanline{8}{μ.prod\ (Measure.infinitePi\ fun\ \_\ :\ ℕ\ ↦\ η.prod\ ν)}
\end{leancode}
\leanlabel{How the Lean proof follows the argument}
\texttt{augment\_\allowbreak{}initial\_\allowbreak{}noise\_\allowbreak{}law} combines the actual product pushforward theorem, the measure-preserving product associator and \texttt{paired\_\allowbreak{}infinite\allowbreak{}Pi}. Almost-everywhere measurable input representatives are transported through the product projection before mapping.
\end{leanbox}

\end{document}
